{\setlength{\parskip}{8mm}
{\huge{\textbf{Resumen}}}

    {\large{El presente Trabajo de Fin de Título (TFT) desarrolla un sistema de diagnóstico de vehículos basado en bus CAN. El sistema implementa los protocolos ISO-TP (ISO 15765-2) como capa de transporte, OBD-II (SAE J1979) y UDS (ISO 14229-1) como capas de aplicación, ejecutándose sobre una Raspberry Pi 4 mediante scripts Python. La sesión de diagnóstico se expone de forma inalámbrica a través de un servidor BLE con perfil GATT, accesible desde una aplicación móvil nativa compatible con Android e iOS que ofrece monitorización en tiempo real, escaneo de DTCs y visualización de logs. El sistema ha sido validado sobre un banco de pruebas basado en un simulador de ECU implementado en un Arduino MKRWAN 1310, tomando como referencia la arquitectura de comunicaciones de los vehículos del grupo VAG (Volkswagen Audi Group), cuya topología de \textit{Gateway} central es representativa de los vehículos modernos de la plataforma PQ25.}}

\newpage

{\huge{\textbf{Abstract}}}

    {\large{This Final Degree Project (TFT) develops a vehicle diagnostics system based on the CAN bus. The system implements ISO-TP (ISO 15765-2) as the transport layer, and OBD-II (SAE J1979) and UDS (ISO 14229-1) as the application layers, running on a Raspberry Pi 4 via Python scripts. The diagnostic session is exposed wirelessly through a BLE server with a GATT profile, accessible from a native mobile application compatible with Android and iOS, offering real-time monitoring, DTC scanning, and log visualization. The system has been validated on a test bench based on an ECU simulator implemented on an Arduino MKRWAN 1310, using the communication architecture of Volkswagen Audi Group (VAG) vehicles as reference, whose central \textit{Gateway} topology is representative of modern PQ25 platform vehicles.}

}}
